\documentclass[a4paper,12pt]{article}
\usepackage[utf8]{inputenc}
\usepackage[T1]{fontenc}
\usepackage[hungarian]{babel}
\usepackage{graphicx}
\usepackage{geometry}
\geometry{a4paper,
		     tmargin = 35mm, 
		     lmargin = 30mm,
		     rmargin = 30mm,
		     bmargin = 30mm}
\usepackage{mathtools}
\usepackage{amsmath}
\usepackage{color}
\usepackage{setspace}
\usepackage{amsmath,amssymb}
\usepackage{float}

\usepackage{indentfirst}

\renewcommand\thesection{\Roman{section}}

\begin{document}

\linespread{1.2}

\begin{titlepage}

	\centering
	\includegraphics[width=0.66\textwidth]{elte.jpg}\par\vspace{1cm}
	{\scshape\LARGE ELTE TTK \par}
	\vspace{3cm}
	{\scshape\Large Az Compton-effektus vizsgálata\par}
	\vspace{1cm}
	{\large\itshape Olar Alex\par}
	\vspace{3cm}
	laborvezető\par
	\vspace{0.3cm}
	{\Large Csanád Máté}
	\vfill
	{\large 2018 \par}
	
\end{titlepage}

\begin{abstract}
\par A mérés célja az volt, hogy a kiértékelés során a lehető legpontosabban járjunk el a számolások közben, hibabecslésen alkalmával. A Compton-effektus tulajdonságait vizsgáltuk. Az energia és a hatáskereszmetszet elméleti szögfüggését vettük alapul a kiértékelés során.
\end{abstract}

\vfill

\tableofcontents

\newpage

\section{Mérési eszközök}

\begin{itemize}
\item plasztik elektron detektor
\item állítható fotodetektor
\item sokcsatornás analizátor (MCA)
\item számítógép
\item $^{137}Cs$, közvetett $\gamma$-foton forrás
\end{itemize}

\vspace{2cm}

\section{Elméleti összefoglaló}

\par A Compton-effektus a fotonok, lazán kötött (kvázi szabad) elektronokon való szóródását írja le, alapegyenlete a négyesimpulzus megmaradásból következik

\begin{equation*}
	E = \frac{E_{0}}{1 + \frac{E_{0}}{m_{e}c^{2}}(1 - cos\vartheta)}	
\end{equation*}

ahol $\vartheta$ az foton szórási szögét, $E_{0}$ annak kezdeti energiáját, míg $m_{e}$ az elektron tömegét jelöli. A mérés során $E_{0}$ értéke adott, a $\gamma$-fotonok $662~keV$ enerigával rendelkeznek. Természetesen $E$ a szórt fotonok energiáját jelöli.
\par Bevezetve a $P = \frac{E}{E_{0}}$ arányszámot, az elméleti differenciális hatáskeresztmetszet felírható

\begin{equation*}
	\frac{d\Sigma}{d\Omega} = \frac{1}{2}r_{0}^{2}(P - P^{2}sin^{2}\vartheta + P^{3})
\end{equation*}

alakban. Ahol $r_{0}$ a klasszikus elektronsugár, melynek értéke a Bohr-modell alapján számolható

\begin{equation*}
	r_{0} = \frac{1}{4\pi\epsilon_{0}}\frac{e^{2}}{m_{e}c^{2}} = 2.82 \times 10^{-15} ~m
\end{equation*}

\par A mérés során ezeket az összefüggéseket próbáljuk a lehető legpontosabban belátni.

\section{Mérési feladatok, kiértékelés}

\par A $^{137}Cs$ minta aktivitása 1963. július 1-jén $486.55~MBq$ volt. A cézium felezési ideje $T = (11018 \pm 9.5)~nap$, míg az adott dátum óta eltelt napok száma a mérés időpontjáig $19 960~nap$ volt. Így az exponenciális bomlástörvény szerint, ma az aktivitás

\begin{equation*}
	A(t) = A_{0}2^{-\frac{t}{T}} = (138.61 \pm 0.08)~MBq
\end{equation*}




\end{document}